% !TEX root = main.tex
\newpage 
\section{Introduction}

\noindent
{\large A reaction-diffusion model for waterborne infection}

\vskip0.2in
We consider a waterborne disease, such as cholera, that can be transmitted 
from both the environment-to-human and human-to-human pathways.
As a starting point, we consider a one-dimensional (1D) spatial domain
$ 0 \leq x \leq a$ for some positive constant $a$
which may represent a river containing the pathogen.

Let the population densities of the susceptible, infected, and recovered individuals be
$S(x, t)$, $I(x, t)$, and $R(x,t)$ respectively, at location $x$ and time $t$.
Let also the pathogen concentration in the river be $B(x,t)$. 
We have the following reaction-diffusion equations to describe 
the disease transmission:
\begin{equation}
\begin{aligned}
 & \frac{\partial S}{\partial t} - c_1 \frac{\partial^2 S}{\partial x^2} 
   = \Lambda - \alpha SB - \beta SI - \mu S, \\
 & \frac{\partial I}{\partial t} - c_2 \frac{\partial^2 I}{\partial x^2} 
   =  \alpha SB + \beta SI - (\mu + \omega + \gamma) I, \\
 & \frac{\partial R}{\partial t} - c_3 \frac{\partial^2 R}{\partial x^2}   
   =  \gamma I-\mu R , \\
 &  \frac{\partial B}{\partial t} + v \frac{\partial B}{\partial x} 
   - d \frac{\partial^2 B}{\partial x^2} = g B \big(1 - \frac{B}{K} \big) +\xi I - \delta B,   
\end{aligned} 
 \label{eqn-recdiff}
\end{equation}
for $ 0 < x < L$ and $ t > 0$.
The parameters $c_1$, $c_2$ and $c_3$ are the diffusion rates of the susceptible,
infected, and recovered individuals,
respectively;
$\Lambda$ is the population influx rate, $\alpha$ is the environment-to-human 
transmission rate,
$\beta$ is the human-to-human transmission rate,
$\mu$ is the natural death rate, $\omega$ is the disease-induced death rate, and
$\gamma$ is the disease recovery rate.
In addition, $g$ is the intrinsic growth rate of the waterborne pathogen,
$K$ is its carrying capacity, $\xi$ is the shedding rate of infected hosts,
$\delta$ is the removal rate of the pathogen from the aquatic environment, 
$d$ is the pathogen diffusion rate, 
and $v$ is the pathogen advection rate that represents the speed of the river flow.


We may use no-flux boundary conditions for system \eqref{eqn-recdiff}:
\begin{equation}
\begin{aligned}
&  \frac{\partial Y}{\partial x}(0, t) = 0, \quad 
  \frac{\partial Y}{\partial x}(L, t) = 0, \quad {\rm where} \;\; Y = S, I, R; \\
& v B(0, t) - d \frac{\partial B}{\partial x}(0, t) = 0 , \quad
 v B(L, t) - d \frac{\partial B}{\partial x}(L, t) = 0, 
\end{aligned}
 \label{bc-noflux}
\end{equation}
for $ t >0$.
Meanwhile, we need to provide appropriate initial conditions:
\begin{equation}
  Y(x,0) = Y_0(x) \geq 0, \quad {\rm where} \;\; Y = S, I, R, B, 
 \label{ic-host}
\end{equation}
for $0 \leq x \leq L$.





